\documentclass[11pt]{article}

\usepackage[margin=1in]{geometry}
\usepackage{setspace}
\usepackage{parskip}
\usepackage{enumitem}
\usepackage{hyperref}
\usepackage{xcolor}
\usepackage{titlesec}
\usepackage{csquotes}

\definecolor{sagpblue}{HTML}{244A5A}
\definecolor{sagpstone}{HTML}{6B6258}

\hypersetup{
    colorlinks=true,
    linkcolor=sagpblue,
    urlcolor=sagpblue,
    citecolor=sagpblue
}

\titleformat{\section}{\Large\bfseries\color{sagpblue}}{}{0pt}{}
\titleformat{\subsection}{\large\bfseries\color{sagpstone}}{}{0pt}{}

\setstretch{1.12}

\title{\textbf{Digital Vision for the Society for Ancient Greek Philosophy}\\
\large A Draft Design Principles Document}
\author{Draft prepared for discussion}
\date{\today}

\begin{document}

\maketitle

\begin{center}
\Large\emph{The Society deserves a digital identity worthy of its intellectual identity.}
\end{center}

\vspace{1em}

\section{Purpose of This Document}

This document is a working draft of the design principles for a renewed digital presence for the Society for Ancient Greek Philosophy. Its purpose is not merely to guide the construction of a website. It is to articulate what the Society's digital identity should express, what virtues it should embody, and how future technical decisions should serve those aims.

The goal is not simply to migrate an existing website from one platform to another. The goal is to create a digital home that faithfully expresses the character, history, intellectual seriousness, and living vitality of the Society.

Technology is not the goal. The goal is to create a place where scholars, students, and visitors immediately sense that they have entered a community devoted to rigorous inquiry, intellectual generosity, and the enduring pursuit of wisdom.

\section{Guiding Principle}

Every design decision should be judged by a single question:

\begin{quote}
Does this help the Society present a digital identity worthy of its intellectual identity?
\end{quote}

If a proposed feature, layout, image, page, or workflow does not serve that purpose, it should be reconsidered.

\section{The Problem We Are Trying to Solve}

The immediate technical problem is that the current website is difficult to maintain and does not fully express the Society's identity. However, the deeper problem is not WordPress, hosting, page templates, or any particular technology.

The deeper problem is this: the Society needs a digital presence that reduces administrative burden while more clearly expressing who the Society is and what it stands for.

The renewed website should help visitors understand the Society quickly, invite participation, preserve institutional memory, and make the work of maintaining the Society easier rather than harder.

\section{Migration, Immigration, and Flourishing}

A simple migration would move the website from one technical platform to another while changing as little as possible. That is not sufficient.

A better metaphor is immigration. The Society's digital presence is moving into a new technological environment, with new expectations, tools, risks, and opportunities. In such a move, identity should be preserved, but growth should be welcomed.

In Aristotelian terms, the aim is not merely preservation but flourishing. The new website should be a better realization of what the Society already is.

\section{Institutional Virtues}

The website should express the virtues of the Society. These virtues are not decorative. They should guide visual design, information architecture, editorial decisions, and technical implementation.

\subsection{Scholarly Excellence}

The Society is devoted to serious scholarship in Ancient Greek philosophy. The website should communicate intellectual depth, rigor, and respect for careful inquiry. It should never feel commercial, superficial, or gimmicky.

\subsection{Intellectual Generosity}

The Society should appear not merely as a gatekeeper of scholarship, but as a community that cultivates scholarship. The site should make clear that discussion, mentorship, feedback, and participation are central to the Society's work.

\subsection{Community}

SAGP is not simply an archive of papers or announcements. It is a scholarly community. Conferences, lectures, works-in-progress sessions, and membership should be presented as forms of shared intellectual life.

\subsection{Continuity}

The Society has a history. The website should honor that history without becoming antiquarian. Tradition should appear as a living inheritance, not as a static monument.

\subsection{Openness}

The Society should feel welcoming to scholars at different career stages, from graduate students to senior scholars, and to participants from a broad international community.

\subsection{Humility}

The digital tone should be quietly confident rather than self-promotional. The Society's seriousness should be evident through clarity, care, and substance, not through exaggerated claims.

\subsection{Beauty}

Beauty is not superficial in this context. A society devoted to ancient philosophy should have a digital presence marked by proportion, restraint, clarity, and grace.

\subsection{Hospitality}

A first-time visitor should not feel lost. A graduate student, prospective member, conference participant, or curious scholar should be able to understand where they are, what the Society does, and how to participate.

\subsection{Living Scholarship}

The website should not feel like an abandoned archive. It should communicate that Ancient Greek philosophy remains an active field of inquiry and that the Society is a living scholarly community.

\section{Desired Visitor Experience}

Within thirty seconds, a visitor should sense that the Society is serious, established, intellectually rich, and welcoming.

Within three minutes, a visitor should understand what the Society does, how to learn more, and how to participate.

Within thirty minutes, a visitor should be able to explore the Society's history, conferences, lectures, membership, and current activities without confusion or friction.

\section{Primary Audiences}

The website should serve several overlapping audiences:

\begin{itemize}[leftmargin=2em]
    \item Current members of the Society.
    \item Prospective members.
    \item Graduate students and early-career scholars.
    \item Senior scholars in Ancient Greek philosophy.
    \item Philosophers in adjacent fields.
    \item Conference participants.
    \item Visitors seeking information about the Society's history, leadership, or activities.
\end{itemize}

The site should not assume that every visitor already knows what SAGP is or how to navigate academic society structures.

\section{Design Language}

The visual language should be timeless rather than trendy, elegant rather than flashy, and modern without sacrificing the Society's identity.

\subsection{Tone}

The tone should be scholarly, calm, generous, and clear. It should avoid both corporate polish and amateur clutter.

\subsection{Typography}

Typography should suggest seriousness, readability, and intellectual tradition. Headings may use strong, elegant forms; body text should be comfortable for sustained reading.

\subsection{Color}

The color palette should be restrained. Stone, parchment, charcoal, deep blue, muted gold, and related tones may be appropriate. Colors should support the content rather than dominate it.

\subsection{Imagery}

Classical imagery may be used, but it should be handled with restraint. Images should evoke intellectual tradition and continuity without becoming kitsch or decoration for its own sake.

\subsection{Whitespace}

Whitespace should be treated as a mark of respect for the content. The site should not feel crowded. Important ideas should have room to breathe.

\subsection{Navigation}

Navigation should be simple, stable, and semantically meaningful. Visitors should be able to find the Society's identity, events, membership information, archives, and contact information easily.

\section{Information Architecture}

A preliminary structure might include:

\begin{itemize}[leftmargin=2em]
    \item Home
    \item About the Society
    \begin{itemize}
        \item Mission and History
        \item Officers and Board
        \item Bylaws
    \end{itemize}
    \item Conferences
    \begin{itemize}
        \item Current Annual Conference
        \item Conference Archive
    \end{itemize}
    \item Distinguished Lectureship
    \item Works in Progress Series
    \item Awards
    \item Membership
    \item Contact
\end{itemize}

This structure should remain flexible until the content has been fully reviewed.

\section{Technology Principles}

Technology should serve the Society's work rather than create new burdens.

\begin{itemize}[leftmargin=2em]
    \item The site should be easy to maintain.
    \item Routine edits should not require specialized technical knowledge.
    \item The system should reduce cognitive load for Society officers.
    \item The site should be secure, stable, and low-maintenance.
    \item The domain and public identity of the Society should remain stable.
    \item The architecture should allow future growth without forcing unnecessary complexity.
\end{itemize}

\section{AI as Design Partner}

AI may be useful not because it replaces human judgment, but because it can help explore the semantic space of possible designs. AI-generated mockups can reveal visual directions that express the Society's identity more clearly than the current site.

However, AI should not invent the Society's content, history, officers, or activities. Authentic content must come from the Society. AI may help organize, draft, summarize, or visually explore, but human stewardship remains central.

\section{What We Will Not Do}

\begin{itemize}[leftmargin=2em]
    \item We will not make technology the center of the project.
    \item We will not sacrifice clarity for novelty.
    \item We will not chase visual trends that will age quickly.
    \item We will not clutter the site with unnecessary features.
    \item We will not allow automated tools to invent institutional facts.
    \item We will not build something that increases Betsy's workload.
    \item We will not confuse copying the old website with preserving the Society's identity.
\end{itemize}

\section{Roadmap}

\subsection{Phase 0: Preservation}

The current website should be archived fully before major changes are made. This includes public pages, media files, PDFs, screenshots, and any available WordPress exports.

\subsection{Phase 1: Vision and Virtues}

Develop and refine this document with Betsy and other stakeholders. Confirm the Society's institutional virtues and desired digital identity.

\subsection{Phase 2: One Page Prototype}

Build one page that expresses the design language. This page should be treated as the experimental prototype: first analyze one page, then build a billion pages.

\subsection{Phase 3: Design System}

Extract reusable patterns from the prototype: typography, headings, page sections, navigation, cards, event layouts, archive layouts, and calls to action.

\subsection{Phase 4: Content Curation}

Review the existing website content. Preserve authentic content, update outdated information, remove duplication, and identify missing material.

\subsection{Phase 5: Full Website Build}

Build the full website using the established design system and curated content.

\subsection{Phase 6: Review and Launch}

Review with Betsy, revise based on feedback, test on desktop and mobile, confirm accessibility basics, and then migrate the domain when ready.

\section{Measures of Success}

The project succeeds if Betsy and the Society's officers can say:

\begin{quote}
This finally feels like the Society.
\end{quote}

Additional measures of success include:

\begin{itemize}[leftmargin=2em]
    \item The site is easier to maintain than the current system.
    \item Visitors can quickly understand what SAGP is and how to participate.
    \item The site feels modern without feeling trendy.
    \item The Society's history and ongoing activities are both visible.
    \item Future officers can understand and preserve the design principles.
\end{itemize}

\section{Closing Statement}

The Society for Ancient Greek Philosophy is devoted to ideas that have endured for millennia. Its digital presence should not merely announce events or store information. It should offer a worthy home for a living intellectual community.

The work ahead is therefore not simply website development. It is digital stewardship.

\end{document}
